\documentclass[11pt]{article}

% ---------- Packages ----------
\usepackage[utf8]{inputenc}
\usepackage[T1]{fontenc}
\usepackage{lmodern}
\usepackage[a4paper,margin=1in]{geometry}
\usepackage{microtype}
\emergencystretch=1.5em

\usepackage{graphicx}
\usepackage{array}
\usepackage{tabularx}
\usepackage{booktabs}
\usepackage{xcolor}
\usepackage{enumitem}
\usepackage{siunitx}
\usepackage{titlesec}
\usepackage{fancyhdr}
\usepackage[hidelinks]{hyperref}
\usepackage[most]{tcolorbox} % callout boxes
\usepackage{float}            % [H] specifier to force inline placement
\usepackage[section]{placeins}% prevent floats moving past section

% ---------- Setup ----------
\setlist{itemsep=0.2em, topsep=0.2em}
\setlength{\parskip}{0.5em}
\setlength{\parindent}{0pt}
\setcounter{secnumdepth}{3}
\setcounter{tocdepth}{2}

\newcolumntype{Y}{>{\raggedright\arraybackslash}X}

% Safe arrows for headings/bookmarks
\newcommand{\gkto}{\texorpdfstring{$\rightarrow$}{->}}
\newcommand{\gkmenu}{(three dots menu)}

% Product/version macros (no content changes—just reusing your values)
\newcommand{\product}{Ghost Key}
\newcommand{\softver}{Software v0.9.2}
\newcommand{\revver}{Revision v0.4}
\newcommand{\docdate}{September 2025}
\newcommand{\docauthor}{Zayn Khan}

% ---------- Callout styles ----------
\newtcolorbox{gkimportant}{
  breakable,
  enhanced,
  colback=orange!5!white,
  colframe=orange!80!black,
  title=Important,
  fonttitle=\bfseries,
  leftrule=1pt, rightrule=1pt, toprule=1pt, bottomrule=1pt,
  arc=2mm, boxrule=0.5pt
}
\newtcolorbox{gknote}{
  breakable,
  enhanced,
  colback=blue!4!white,
  colframe=blue!70!black,
  title=Note,
  fonttitle=\bfseries,
  leftrule=1pt, rightrule=1pt, toprule=1pt, bottomrule=1pt,
  arc=2mm, boxrule=0.5pt
}

% ---------- Header / Footer ----------
\pagestyle{fancy}
\fancyhf{}
\lhead{\product{} — Technical Documentation}
\rhead{\softver{} \ (\revver)}
\cfoot{\thepage}

% ---------- Title ----------
\title{\product: Technical Documentation}
\date{\docdate}

\begin{document}
\maketitle

% ----- Document control (standard manual front-matter) -----
\section*{Document Control}
\begin{tabularx}{\linewidth}{@{}lY@{}}
\toprule
\textbf{Document Title} & \product: Technical Documentation \\
\textbf{Version}        & \softver{} \ (\revver) \\
\textbf{Date}           & \docdate \\
\textbf{Owner}          & \docauthor \\
\textbf{Status}         & Draft for review \\
\bottomrule
\end{tabularx}

% ----- TOC / LOF -----
\tableofcontents
\listoffigures
\clearpage

% =====================================================================
%                      TECHNICAL DOCUMENTATION CONTENT
% =====================================================================

\textbf{\Large Ghost Key --- Technical Documentation}

\emph{\revver{} --- for review}\\
\emph{\product{} Software \softver}

\medskip\hrule\medskip

\section{At a glance}
\begin{itemize}
  \item \textbf{What it is:} ESP32-based vehicle access \& start system with RFID + Bluetooth authentication, push-to-start, and a web configuration portal.
  \item \textbf{Core flows:} Authentication, accessory/ignition/start, immobilizer control.
  \item \textbf{Interfaces:} Start button, brake input, system LED, button LED, buzzer, RFID reader, BLE, and web UI.
\end{itemize}

\section{Quick start}
\begin{enumerate}
  \item \textbf{Power on.} System LED blinks \(\sim\) once per second.
  \item \textbf{Enter Config Mode.} Hold \textbf{Start} for \textbf{10s} (brake \textbf{not} pressed). The \textbf{button LED pulses smoothly}.
  \item \textbf{Join Wi-Fi.} Connect to \textbf{Ghost Key Configuration} (default password \textbf{123456789}).
  \item \textbf{Open the portal.} Browse to \url{http://192.168.4.1} in \textbf{Safari (iOS)} or \textbf{Chrome (Android)}.
  \item \textbf{Set up.} Choose \textbf{Ghost Key}, \textbf{Ghost Power}, or \textbf{Both}. Change the \textbf{Wi-Fi password} and set an \textbf{Admin PIN} for the web portal.
  \item Pair \textbf{Bluetooth} and enroll any additional \textbf{RFID} (up to 5).
  \item \textbf{Calibrate Proximity Detection.} Run \textbf{Bluetooth Calibration} after installation for best proximity accuracy (see \emph{Authentication} \gkto{} \emph{Bluetooth}).
  \item Set \textbf{Starter Pulse} and \textbf{Auto-lock}.
  \item \textbf{Exit Config.} Press \textbf{Start} once. Test \textbf{lock/unlock} and \textbf{start/stop}.
\end{enumerate}

\begin{gkimportant}
\textbf{Important:} Setup \emph{must} be completed even if settings are not changed. Device will drain the car battery until setup is completed.\\
\textbf{Important:} If Ghost Key \emph{and} Bluetooth are enabled, a Bluetooth device \textbf{must} be paired; otherwise the device will not come out of deep sleep.
\end{gkimportant}

\section{Daily use}
\subsection{Start \& stop}
\begin{enumerate}
  \item \textbf{Authenticate} via RFID or a paired phone (Bluetooth).
  \item \textbf{Press and hold the brake.}
  \item \textbf{Press Start} to crank; release after the engine starts.
  \item \textbf{Press Start (with brake)} to stop the engine.
\end{enumerate}

\subsection{Accessory / ignition without starting}
\begin{itemize}
  \item \textbf{Press Start (no brake)} to cycle \textbf{OFF \gkto{} ACC \gkto{} IGN \gkto{} OFF}.
\end{itemize}

\subsection{Security \& auto-lock}
\begin{itemize}
  \item \textbf{Immobilizer/Security relay} engages when the engine is off and there is \textbf{no valid authentication}.
  \item System \textbf{auto-locks} by default \textbf{30 s} after shutdown (configurable \textbf{5--120 s}).
  \item Any valid authentication \textbf{unlocks/disarms}.
\end{itemize}

\section{Controls \& indicators}
\subsection{Start button (actions)}
\begin{itemize}
  \item \textbf{Press (no brake):} Cycle \textbf{OFF \gkto{} ACC \gkto{} IGN \gkto{} OFF}
  \item \textbf{Press with brake:} Start engine (if authenticated) or stop when running
  \item \textbf{Press \& hold 10s (no brake):} Enter/exit \textbf{Config Mode} (button LED pulses smoothly)
  \item \textbf{Press \& hold 30s (with brake):} \textbf{Factory reset} (settings only; RFID keys preserved)
\end{itemize}

\subsection{LEDs}
\textbf{System LED} (status/heartbeat)

\begin{tabularx}{\linewidth}{@{}lY@{}}
\toprule
\textbf{Pattern} & \textbf{Meaning}\\
\midrule
Blink every \(\sim\)1s & Active power mode \\
Blink every \(\sim\)2s & Light sleep \\
Off & Deep sleep or powered off \\
\textbf{3 quick flashes} & Entering \textbf{Config Mode} \\
\textbf{Goes off} & Exiting \textbf{Config Mode} \\
\textbf{Rapid \(\times\)10} & Factory reset in progress \\
\textbf{3 quick 50 ms flashes} & (see Button LED) \\
\bottomrule
\end{tabularx}

\medskip
\textbf{Start Button LED} (illumination/feedback)

\begin{tabularx}{\linewidth}{@{}lY@{}}
\toprule
\textbf{State} & \textbf{Meaning}\\
\midrule
\textbf{Solid} & Brake pressed \textbf{and} authenticated \\
\textbf{Off} & Not authenticated \textbf{or} brake not pressed \\
\textbf{Smooth pulse} & \textbf{Config Mode} active \\
\textbf{3 quick flashes} & Authentication error / access denied \\
\bottomrule
\end{tabularx}

\subsection{Buzzer}
\begin{tabularx}{\linewidth}{@{}lY@{}}
\toprule
\textbf{Sound} & \textbf{Meaning}\\
\midrule
\textbf{1 short beep} & RFID authenticated (successful read) \\
\textbf{2 short beeps} & Unauthorized RFID tag \\
\textbf{Continuous 2 s} & Entering \textbf{Config Mode} (Ghost Power only) \\
\textbf{Continuous 1 s} & Exiting \textbf{Config Mode} (Ghost Power only) \\
\bottomrule
\end{tabularx}

\section{Authentication}
\subsection{RFID}
\begin{itemize}
  \item \textbf{Capacity:} Store up to \textbf{5} tags
  \item \textbf{Validation:} Requires \textbf{one} good read
  \item \textbf{Timeout:} Auth remains valid for \(\sim\)\textbf{30 s} after tag removal
  \item \textbf{Range:} Close proximity (a few cm)
\end{itemize}

\subsection{Bluetooth (BLE)}
\begin{itemize}
  \item \textbf{Device name:} \textbf{Ghost Key}
  \item \textbf{Max bonds:} \textbf{3} devices
  \item \textbf{Method:} RSSI-based \textbf{confidence} with hysteresis
  \begin{itemize}
    \item \textbf{Wake thresholds (during BLE scan window):}
    \begin{itemize}
      \item \(\le\) \textbf{35\%} \gkto{} remain in \textbf{Deep sleep} (continue 10s off / 10s scan cycling)
      \item \textbf{35--45\%} \gkto{} enter \textbf{Light sleep} at end of window (quicker response)
      \item \(\ge\) \textbf{45\%} \gkto{} wake to \textbf{Active}
    \end{itemize}
    \item \textbf{Auth thresholds (permission to start):}
    \begin{itemize}
      \item \textbf{Grant} at \(\ge\) \textbf{65\%} confidence
      \item \textbf{Revoke} at \(\le\) \textbf{50\%} confidence
    \end{itemize}
  \end{itemize}
  \item \textbf{Range:} \(\sim\)3--10 ft depending on phone and environment
  \item \textbf{Calibration (recommended):} Place your phone at the desired ``OK to start'' location (e.g., in your pocket in the \textbf{driver’s} seat) and run the \(\sim\)\textbf{30 s} calibration.
\end{itemize}

\begin{gknote}
\textbf{Note:} In \textbf{Deep sleep}, RFID is \textbf{off}. Bluetooth is scanned during the 10s window to evaluate confidence and wake decisions. If Bluetooth is disabled the start button \emph{or} the brake pedal can be used as a wakeup to use RFID.
\end{gknote}

\section{Configuration portal}
\subsection{Connect}
\begin{enumerate}
  \item Enter \textbf{Config Mode} (hold \textbf{Start} 10 s, no brake).
  \item Join \textbf{Wi-Fi:} \texttt{Ghost Key Configuration}\\
        (default password \textbf{123456789})
  \item Browse to \url{http://192.168.4.1} and sign in with your \textbf{Interface PIN}.
\end{enumerate}

\begin{gkimportant}
\textbf{Important:} On first setup, change the \textbf{Wi-Fi password} and the \textbf{Interface PIN}.
\end{gkimportant}

\subsection{Add to Home Screen (optional)}
\begin{itemize}
  \item \textbf{iOS / iPadOS (Safari only):} Share \gkto{} \textbf{Add to Home Screen} while on \textbf{192.168.4.1}.
  \item \textbf{Android (Chrome recommended):} Menu \gkmenu{} \gkto{} \textbf{Add to Home Screen} while on \textbf{192.168.4.1}.
\end{itemize}

\subsection{Key settings}
\begin{itemize}
  \item \textbf{Modes:} Enable/disable \textbf{Ghost Key} or \textbf{Ghost Power}, or enable \textbf{Both}
  \item \textbf{RFID keys:} Add, name, or remove (max \textbf{5})
  \item \textbf{Bluetooth:} Pairing mode, device names, remove bonds (max \textbf{3})
  \item \textbf{Starter pulse:} \SI{100}{\milli\second}--\SI{3000}{\milli\second} (default \SI{700}{\milli\second})
  \item \textbf{Auto-lock delay:} \SI{5}{\second}--\SI{120}{\second} (default \SI{30}{\second})
  \item \textbf{Calibration:} Bluetooth proximity calibration (\(\sim\)\textbf{30 s})
  \item \textbf{Developer Diagnostics:} Live RSSI/confidence, status, logs
\end{itemize}

\section{Troubleshooting}
\subsection{Quick checks}
\begin{itemize}
  \item \textbf{No LED:} Check power and fuses; unit may be in \textbf{Deep sleep} (press brake or Start).
  \item \textbf{Rapid LED flashing:} Authentication failed; present a valid RFID or approach with a paired phone.
  \item \textbf{LED stuck on:} Power-cycle; if it persists, reflash firmware or contact support.
  \item \textbf{No buzzer on RFID:} Tag not enrolled or reader wiring/placement issue.
\end{itemize}

\subsection{Symptom \gkto{} fix}
\begin{tabularx}{\linewidth}{@{}lYY@{}}
\toprule
\textbf{Symptom} & \textbf{Likely cause} & \textbf{Try this}\\
\midrule
Won’t start with button & Not authenticated & Present RFID or bring paired phone closer; check BLE is enabled \\
Starts then stalls & Ignition outputs or immobilizer timing & Verify \textbf{IGN1/IGN2} wiring; increase \textbf{Starter Pulse} \\
Can’t enter Config Mode & Brake is pressed or device is not authenticated & Release brake; authenticate with RFID or Bluetooth; hold \textbf{Start} for \textbf{10 s} \\
BLE unreliable & Phone/environment variability & Re-run \textbf{Calibration}; relocate antenna \\
RFID intermittent & Tag orientation or interference & Reposition tag/antenna; clean tag; ensure \textbf{two consecutive} reads \\
\bottomrule
\end{tabularx}

\section{Reset \& recovery}
\begin{itemize}
  \item \textbf{Exit Config Mode:} Press \textbf{Start} once (or \textbf{hold an enrolled RFID for 10s} in \textbf{Ghost Power-only} systems).
  \item \textbf{Factory reset:} Hold \textbf{Start + brake}. After 10s the system LED will start to toggle each second. Continue holding for \textbf{30s} to complete the reset.
  \begin{itemize}
    \item \textbf{Preserves:} RFID keys
    \item \textbf{Resets:} All other settings and Bluetooth bonds
    \item \textbf{Feedback:} System LED toggles each second; rapid flashes at completion
  \end{itemize}
\end{itemize}

\section{Specifications (summary)}
\begin{itemize}
  \item \textbf{MCU:} ESP32 (240/160/80 MHz power scaling)
  \item \textbf{Auth:} RFID (EM4095-style), BLE proximity with confidence model
  \item \textbf{Outputs:} ACC, IGN1, IGN2, START, SECURITY (relay-driven)
  \item \textbf{Inputs:} Start button, brake, accessory sense
  \item \textbf{Indicators:} System LED, button LED, buzzer
  \item \textbf{Wi-Fi:} AP configuration (captive portal style)
\end{itemize}

\clearpage
\section*{Appendix A — Installer reference}
\subsection*{GPIO map}
\begin{tabularx}{\linewidth}{@{}l l@{}}
\toprule
\textbf{Function} & \textbf{GPIO}\\
\midrule
System LED           & 23 \\
Button LED           & 26 \\
Beeper               & 25 \\
Start button (input) & 35 (active-low, ext pull-up) \\
Brake input          & 18 (active-low) \\
Accessory input      & 19 (active-low) \\
RFID Demod Out       & 12 \\
RFID SHD             & 13 \\
RFID MOD             & 27 \\
RFID RDY/CLK         & 14 \\
ACC relay            & 15 \\
IGN1 relay           & 5 \\
IGN2 relay           & 4 \\
START relay          & 16 \\
SECURITY relay       & 17 \\
\bottomrule
\end{tabularx}

\subsection*{Power \& sleep behavior (detailed)}
\begin{itemize}
  \item \textbf{Deep sleep cycle:} \textbf{4s radio-off / 8s BLE-scan} cycling; \textbf{RFID OFF}. During each BLE window: \(\le\) \textbf{35\%} confidence \gkto{} remain in \textbf{Deep sleep}; \textbf{35--45\%} \gkto{} transition to \textbf{Light sleep} at the end of the window; \(\ge\) \textbf{45\%} \gkto{} wake to \textbf{Active}.
  \item \textbf{Light sleep:} Borderline-confidence mode (\textbf{35--45\%}) for quicker response; BLE remains active.
  \item \textbf{Auth hysteresis:} \textbf{Grant} at \(\ge\) \textbf{65\%} confidence; \textbf{Revoke} at \(\le\) \textbf{50\%}.
  \item \textbf{Wake stabilization:} \(\sim\)\textbf{2 s} after deep-sleep wake.
  \item \textbf{User wake:} \textbf{Start} always wakes; device then authenticates via BLE or reads a paired RFID fob.
  \item See Figure 1 below for an overview
\end{itemize}

\subsubsection*{Power State Diagram}
\IfFileExists{Mermaid Chart Diagrams Aug 29 2025 (1).png}{%
\begin{figure}[H]
  \centering
  \includegraphics[width=\linewidth]{Mermaid Chart Diagrams Aug 29 2025 (1).png}
  \caption{Power State Diagram}
\end{figure}
}{%
\begin{figure}[H]\centering
\fbox{\parbox{\linewidth}{\centering [Power State Diagram image placeholder]}}
\caption{Power State Diagram}
\end{figure}
}

\subsection*{Authentication Logic Flow}
\IfFileExists{Authentication Logic from Google Drive.png}{%
\begin{figure}[H]
  \centering
  \includegraphics[width=\linewidth]{Authentication Logic from Google Drive.png}
  \caption{Authentication Logic Flow}
\end{figure}
}{%
\begin{figure}[H]\centering
\fbox{\parbox{\linewidth}{\centering [Authentication Logic Flow image placeholder]}}
\caption{Authentication Logic Flow}
\end{figure}
}

\subsection*{State Diagram}
\IfFileExists{Ghost Key State Machine.png}{%
\begin{figure}[H]
  \centering
  \includegraphics[width=\linewidth]{Ghost Key State Machine.png}
  \caption{State Diagram}
\end{figure}
}{%
\begin{figure}[H]\centering
\fbox{\parbox{\linewidth}{\centering [Authentication State Diagram image placeholder]}}
\caption{Authentication State Diagram}
\end{figure}
}

\section*{Change log}
\begin{itemize}
  \item v0.4: Added State Machine Diagram
  \item v0.3.1: Slight edit to quick start to include note about completing setup
  \item \textbf{v0.3}: Added Power State Diagram
  \item \textbf{v0.2:} Reorganized sections; unified wake vs.\ auth thresholds; corrected IP/browser notes; standardized terms (Wi-Fi, Admin PIN); consolidated calibration guidance; clarified LED/buzzer behaviors; moved detailed sleep logic to Appendix.
  \item \textbf{v0.1:} Initial structure + content from codebase analysis
\end{itemize}

\end{document}
